\dftfe{} is a C++ code for materials modeling from first principles using Kohn-Sham density functional theory.
It is based on adaptive finite-element discretization that handles all-electron and pseudopotential calculations in the 
same framework, and incorporates scalable and efficient solvers for the solution of the Kohn-Sham equations. Importantly, \dftfe{} 
can handle general geometries and boundary conditions, including periodic, semi-periodic and non-periodic systems. \dftfe{} code 
builds on top of the deal.II library for everything that has to do with finite elements, geometries, meshes, etc., and, through 
deal.II on p4est for parallel adaptive mesh handling.

\subsection{Authors}
\label{sec:authors}
\dftfe{} is developed and maintained by the \href{http://www-personal.umich.edu/~vikramg/}{Computational Materials Physics
group at University of Michigan} (directed by Prof. Vikram Gavini, Dr. Sambit Das), and the \href{https://sites.google.com/view/matrix-lab}{MATRIX lab, Indian Institute of Science} (directed by Prof. Phani Motamarri). \dftfe{} code is currently maintained by a group of development leads/mentors and  principal developers, 
who manage the architecture of the code and the core functionalities. Finally, all contributors who have contributed to major parts
of the \dftfe{} code or sent important fixes and enhancements are listed under contributors. All the underlying 
lists are in alphabetical order based on the last name.

\paragraph{Mentors/Development Leads}
\begin{itemize}
    \item Sambit Das (University of Michigan, USA 2017-present)    
     \item Vikram Gavini (University of Michigan, USA 2015-present)
     \item Phani Motamarri (University of Michigan, USA 2015-2019; Indian Institute of Science, India 2019-present)
\end{itemize}

\paragraph{Principal developers}
\begin{itemize}
	\item Sambit Das (University of Michigan, USA 2017-present)
	\item Krishnendu Ghosh (University of Michigan, USA 2017-2020)    
	\item Phani Motamarri (University of Michigan, USA 2015-2019; Indian Institute of Science, India 2019-present)
	\item Shiva Rudraraju (University of Michigan, USA 2015-2018)
 \item Nikhil Kodali (Indian Institute of Science, India 2021-present)
\item Kartick Ramakrishnan (Indian Institute of Science, India 2020-present)   
\end{itemize}



\paragraph{Contributors}
\begin{itemize}
    \item Sambit Das (University of Michigan, USA 2017-present)
	\item Krishnendu Ghosh (University of Michigan, USA 2017-2020)   
 	\item Phani Motamarri (University of Michigan, USA 2015-2019; Indian Institute of Science, India 2019-present) 
	\item Shiva Rudraraju (University of Michigan, USA 2015-2018)    
 \item Nikhil Kodali (Indian Institute of Science, India 2021-present)
\item Kartick Ramakrishnan (Indian Institute of Science, India 2020-present) 
	\item Bikash Kanungo (University of Michigan, USA 2023-present)
 \item Shukan Parekh (University of Michigan, USA Summer 2019)
 \item Rudra Panch (Indian Institute of Technology, Madras, India 2024-present)
  \item Ian C. Lin (University of Michigan, USA 2018-2019)
\item Srinibas Nandi (Indian Institute of Science, India 2023-present) 
\item Gourab Panigrahi (Indian Institute of Science, India 2021-present)
  	\item Vishal Subramanian (University of Michigan, USA 2022-present)
\item David Rogers (Oak Ridge National Laboratory, USA 2020-2023)
\item Denis Davydov (University of Erlangen-Nuremberg, Germany 2018)
	
\end{itemize}

\subsection{Acknowledgments}
The development of \dftfe{} open source code relating to pseudopotential calculations has been funded in part by the Department 
of Energy PRISMS Software Center at University of Michigan, and the Toyota Research Institute. The development of \dftfe{} open 
source code relating to all-electron calculations has been funded by the Department of Energy Basic Energy Sciences. The methods 
and algorithms that have been implemented in \dftfe{} are outputs from research activities over many years that have been 
supported by the Army Research Office, Air Force Office of Scientific Research, Department of Energy Basic Energy Sciences, 
National Science Foundation and Toyota Research Institute. 

Part of the DFT-FE open source code development related to the capabilities --- PAW, Ab initio MD, NEB, non-collinear magnetism with spin-orbit coupling has been funded by seed grant from Indian Institute of Science; Science and Engineering Research Board (SERB) Startup Research Grant (SRG); National Supercomputing Mission (NSM) R\&D for Exascale Grant from the Department of Science and Technology (DST), India; Indo-Korean Science and Technology Center (IKST), Bangalore [South Korea] and Google India Research Award.

\subsection{Referencing \dftfe{}}
Please refer to \href{https://sites.google.com/umich.edu/dftfe/referencing}{referencing  \dftfe{}} to properly cite the use of 
\dftfe{} in your scientific work. 
